\begingroup
\scriptsize
\setlength{\tabcolsep}{4pt}
\begin{tabularx}{\textwidth}{c X c c c c c l}
\toprule
\# & Configuration & MTMC IDF1 & $\Delta$ (pp) & IDF1 & MOTA & HOTA & Source \\
\midrule
1 & Baseline CC (basic PCA 384D, no FIC, no intra-merge) & 74.0\% & --- & --- & --- & --- & 10c v14--v22 grid \\
2 & + FIC whitening (reg=0.1 in the historical ablation notes) & 76.0\% & +2.00 & --- & --- & --- & 10c v31--v33 \\
3 & + Power normalization ($\alpha=0.5$) & 76.5\% & +0.50 & --- & --- & --- & 10c v28--v29 \\
4 & + AQE ($K=3$, $\alpha=5.0$) & 77.4\% & +0.90 & --- & --- & --- & 10c v14 vs. v22 \\
5 & + Conflict-free CC algorithm & 77.61\% & +0.21 & --- & --- & --- & 10c v23--v26 \\
6 & + Intra-camera merge (thresh=0.80, gap=30) & 78.28\% & +0.67 & --- & --- & --- & 10c v34 \\
7 & + min\_hits=2 (historical best) & 78.4\% & +0.12 & 79.8\% & --- & --- & 10c v44 / historical v80 recipe \\
8 & Current reproducible best (v80-restored recipe on current codebase) & 77.5\% & --- & 80.1\% & 67.1\% & 58.1\% & 10c v52 \\
\bottomrule
\end{tabularx}

\vspace{0.4em}
\par\noindent\footnotesize Controlled ON/OFF ablations also show a +0.9pp gain from the temporal-overlap bonus (10c v75), but that result was measured on a different branch, so it is reported in the text rather than inserted into this strictly cumulative sequence. Unavailable intermediate IDF1, MOTA, and HOTA values are shown as em dashes rather than inferred.
\endgroup